\scanpage{38}

(et non propre), en utilisant un \unsure{relèvement} de $X_0$, non comme schéma ni comme schéma formel, mais en quelque chose qui ressemble à une \uline{variété rigide-analytique} à la Tate. De ce fait, il n'est défini, \inserted{dans $[\,]$} que des invariants \uline{locaux}. Par communication orale de W., il serait arrivé à définir des invariants globaux (sur $W$, et non seulement sur $K$), en travaillant \inserted{avec le complexe de De Rham} sur le \uline{site des relèvements de W.M} \deleted{(\illegible{3})} \uline{des ouverts affines}, et sauf erreur à justifier les assertions que \unsure{j'ai faites}.

c) On voit, grâce au th. de comparaison, que par cette construction (quand $X_0$ \unsure{se} relève en $X$ lisse sur $S$) on trouve des $H^i(X_0)$ qui donnent les bons nbs de Betti --- ce sont ceux de la fibre générique $X_t$ !

\margin{37}
